\documentclass[11pt]{article}
\usepackage[margin=25mm]{geometry}\usepackage[T1]{fontenc}\usepackage{lmodern}
\usepackage{amsmath,amssymb,booktabs,hyperref}
\setlength{\emergencystretch}{3em}
\title{Response to scientific review}
\author{Gregory Morse \and Peter Rakyta}\date{7 November 2026}
\begin{document}\maketitle
We thank the reviewer for independently checking the identities, finite-alphabet
search, invariant inequality, and lower-bound argument. The revision follows
the theory-paper direction. Its title is now \emph{CNOT costs of controlled
balanced involutions and an invariant two-SQiSW criterion}. The atlas is an
accompanying reference resource. All seven numbered concerns have prompted
changes to the text, data presentation, experiments, or claim scope.

\section*{1. Grid size, obstruction types, and optimality denominators}
We agree that the headline grid size hid the distribution of difficulty.
The manuscript now reports the named targets and fixtures separately and
gives one complete table per phase policy, broken down by physical width.
The 217 named one-qubit entries, 192 phased Cliffords, and standard determinant
and discrete-alphabet obstructions are explicitly identified. Zero-entangler
constructions are separated from positive-entangler constructions.

A new row-resolved CSV reconciles every retained cell with its explicit
matched proof sources. We take a union rather than adding overlapping audits;
fixed-target minima and minimum uniform family budgets are disjoint columns.
Legacy Boolean optimality flags alone are not counted as certificates.
Among 2,295 positive-entangler literal constructions, the union contains
732 fixed-target minima and 710 uniform family minima; 853 remain unproved
under this accounting. Among 3,057 effective positive-entangler projective
constructions, it contains 984 fixed-target minima and 955 uniform family
minima; 1,118 remain unproved. These are entangler certificates, not local-cost
optimality claims.

The reported 1,916 projective denominator counts explicitly stored projective
results. A further 1,141 positive-entangler literal circuits were reused
projectively, giving the complete-grid denominator 3,057. The old raw flags
694/423 and selected-audit counts were different sets and were inadequately
explained; neither is presented as the new union. The new CSV records both
the raw flag and the matched proof sources so this distinction can be checked.

The obstruction table identifies all 2,601 literal determinant-group cases.
All 268 projective obstructions occur in the three discrete presets. These
are described as standard feasibility tests, with exceptional parameter
specializations distinguished from whole-family representation.

\section*{2. Clifford reference and missing atlases}
We removed any framing of the 192 one-qubit targets as a major unexplored
synthesis class. Matsumoto--Amano, Kliuchnikov--Maslov--Mosca, and
Giles--Selinger are now cited. The finite-alphabet shortest-path result is
presented as an explicit reference and regression certificate, with no
continuous-alphabet global-local-optimality extrapolation. We did not add a
two-qubit Clifford classification or imply that it has been completed.

The complete gate and Clifford PDFs and their editable sources are now
included in the full review ZIP, rather than only pointed to elsewhere.
The manuscript's coverage tables carry the main quantitative evidence;
atlas page count is not used as evidence of algorithmic novelty.

\section*{3. Prior art, generalization, and proof detail}
We added Yu--Duan--Ying, Yu--Ying, Shende--Markov, and Nielsen et al., and
explain how their lower-bound and controlled-product arguments relate to
ours. Ordinary Fredkin, Toffoli, and controlled fermionic SWAP have different
relative blocks and resource counts. Unrestricted two-qubit gates are kept
separate from Ising-class interactions.

The main new generalization classifies every controlled balanced two-qubit
involution: $V=V^\dagger$, $V^2=I$, $\operatorname{tr}V=0$. A Pauli-coefficient
normal form proves costs four at OSR four, three at OSR two, and one or two
for local products according to support. It includes the zero-vector and
rank-deficient degeneracies and gives exact symbolic boundary constructions.
There is no rank-three stratum. We separately prove that nonlocal balanced
relative blocks need three unrestricted two-qubit gates and provide a
three-gate upper construction. The fermionic SWAP example therefore costs
four Ising interactions but only three unrestricted interactions.

We do not assert the suggested classification of arbitrary controlled-$V$
by $V$'s Weyl class alone: $I$, $Z_a$, and $Z_aZ_b$ have the same local Weyl
class but controlled costs zero, one, and two. The corrected generalization
uses the additional involution and support structure explicitly.

The four-Ising proof now defines control OSR, derives the target cut ranks,
counts all twelve surviving interaction words, and expands the conjugation
bound into $M_0\otimes I+M_1\otimes Z$. Fermionic SWAP's Pauli expansion
shows OSR four directly and its local equivalence to iSWAP is stated.
The improved literal witness has four CNOTs and six Ry rotations, replacing
the former ten-local-gate witness. Exact tests cover the free two-angle
normal form and every rank-two boundary, not only numerical samples.

\section*{4. Baselines and application comparison}
We executed Qiskit 2.6.2, pytket 2.18.5, and BQSKit 1.2.1 on eight independent
fixed targets, using both library-source and matrix routes where supported.
There are 42 SDK attempts. We additionally executed Quartz and QUESO on the
eight source circuits each, preserving raw outputs, revisions, equivalence
rules, harness changes, budgets, and fixed-wire full-matrix replay. Their
objective is total NAM gates; only one-qubit fusion changes their reported
local alphabet. UniversalQCompiler remains unrun because its licensed kernel
is unavailable; that explanation is not applied to the open-source tools.

BQSKit matches four CNOTs and six locals for fFredkin. Our separate 108-layout
BQSKit experiment finds a four-CNOT fit in about six seconds. These numerical
findings support the reviewer's observation that the upper bound is easy;
the lower proof and generalized classification carry the mathematical
contribution. All eleven older SDK entangler wins are now explicitly
identified as fFredkin across different presets. Compiler losses are also
reported, including Quartz's verified 30-CNOT C4X result versus the current
36-CNOT atlas entry, and smaller C3X/SSWAP local counts from other tools.
No claim of universal dominance remains.

The tomography-figure 27-to-3 comparison is removed from the active manuscript
and evidence appendix. The latter's incomplete reference is replaced with
Warren et al.'s full title, authors, venue, year, and DOI. Published C3X/C4X
and ordinary Fredkin constructions are discussed with matching ancilla and
resource assumptions and are identified as upper bounds, not automatically
global minima.

The parameter-reuse table now prominently includes independent numerical
recompilation: all ten bound instances give 18 CNOTs, 40 locals, depth 25,
at a mean of about 7.68 ms. This beats the symbolic 18/42/31 circuit in local
count and depth. The older 36-CX gatewise lowering is explicitly a weak
contextual control, and the known three-CX fSim source is acknowledged.
We withdraw the inference of a new 36-to-18 resource advance over analytical
fSim synthesis. Cold synthesis, SDK compile time, and parameter binding are
separated.

\section*{5. OSR ablation and practical utility}
We added a construction ablation disabling OSR factorization and residual
peeling while retaining native charts, Clifford/Pauli construction, and
posthoc verification. It gives the same two-qubit success count as the full
method, 11/12, with one deadline each. The prior chart-disabled result is
3/12. The added opaque three- and four-qubit product probes remain unresolved
for both methods in this configuration. These outcomes are retained and
reported; they do not demonstrate an OSR-search advantage on this workload.
The title no longer attributes these chart successes to OSR-guided search.
OSR's role in the theorems is established by proof instead.

We also implemented a bounded atlas peephole pass on eight existing benchmark
sources in two native presets. Every proposed replacement requires exact
identity acceptance; the control is the same gatewise native lowering.
The pass finds no improving windows on the suite. A separately labeled
three-CX SWAP positive control checks its exact identity and native count
accounting. We report this negative suite outcome and make no general
peephole speedup claim. This is evidence about the current optimizer's limits,
not hidden or counted as successful application improvement.

\section*{6. Secondary score and symbolic-expression cost}
The parameter-description score and its dependence on simplification are
moved to an appendix. It is not a calibrated pulse model and does not enter
the main entangler theorems. The historical negative-cycle repair is no longer
part of the scientific contribution. Unknown equal-entangler scores remain
incomparable.

The long SQiSW chart is now called a machine-inspectable scalar DAG, not a
short directly readable expression. A new optional compiled numerical
evaluator shares dependencies and preserves lazy degeneracy branches.
For the 585-binding chart, compilation takes about 1.36 s and 1,000 warm
bindings about 0.97 s in the recorded run. Full matrix replay passes at 64
points, including singular corners and angle folds. This checks practical
reuse of an already exactly certified family, without claiming that it always
beats numerical KAK or that sampling proves an identity. The text explains
when symbolic identity and phase contracts are useful and explicitly retains
numerical specialization as a competitive alternative.

\section*{7. Presentation and editorial cleanup}
The manuscript is theory-first, with contributions stated directly and
scope limits consolidated. It names Gregory Morse and Peter Rakyta, removes
anonymous-review labeling, cites every displayed figure in the text, and
expands the control-rank argument. Affiliations and individual human
contribution statements remain for the authors to supply before submission;
none has been invented.

We removed coordination prose from the proof supplement and scientific
archive view. Original historical records remain immutable in the repository;
the current review package uses the cleaned scientific view and current
sources. The new response, evidence appendix, source bundles, complete
atlases, and verification manifests describe the revised claims consistently.
AI assistance is disclosed as assistance, not authorship or proof authority.

\section*{Additional SQiSW result}
We proved the iff statement suggested by the review. For chamber coordinates,
write $g_1=(u+iv)^2$ with $u\ge0$; then $|g_1|+\Re g_1=2u^2$. The two
invariant inequalities are equivalent to $W\ge4u$. Their difference factors
as
\[
 W-4u=2[\cos2x-\cos2(y+|z|)]
          [\cos2x-\cos2(y-|z|)].
\]
The chamber orders fix the sign of the second factor; $z=0$ is handled
separately. This proves equivalence to $x\ge y+|z|$, including the closed
boundary, and Huang et al.'s known region gives necessity and sufficiency.
We credit the region to Huang et al. and claim the invariant reformulation
and checked implementation, not a new pulse-cap theorem.
\end{document}
